%%
%% fall-lecture13.tex
%% 
%% Made by Alex Nelson <pqnelson@gmail.com>
%% Login   <alex@lisp>
%% 
%% Started on  2026-10-27T10:44:32-0700
%% Last update 2026-10-27T10:44:32-0700
%% 

\lecture{}

\begin{recall}
We determined $H^{*}(G_{n};\ZZ_{2})\iso\ZZ_{2}[w_{1}(\gamma^{n}),\dots,w_{n}(\gamma^{n})]$.
We proved uniquess of the Stiefel--Whitney classes (which, we stress,
involved the universal property of $\gamma^{n}$). This raised the
question: how necessary is the universal property for proving
uniqueness? We can ``kind of'' bypass this reliance on the ``general''
universal property of $\gamma^{n}$ by relying on the splitting
principle plus the universal property of $\gamma^{1}$.
\end{recall}

\begin{theorem}[Splitting principle]
Let $\xi\colon E\to X$ be an $\RR^{n}$-bundle, where $X$ is a
paracompact space. Then there exists a topological space $Y$ and a map
$p\colon Y\to X$ such that
\begin{enumerate}
\item $p^{*}\colon H^{*}(X;\ZZ_{2})\to H^{*}(Y;\ZZ_{2})$ is injective
\item the pullback bundle $p^{*}\xi$ splits
  $p^{*}E=L_{1}\oplus\dots\oplus L_{n}$ as the Whitney sum of line
  bundles $L_{1}$, \dots, $L_{n}$ over $Y$.
\end{enumerate}
\end{theorem}

\begin{corollary}
The Stiefel--Whitney class $w(\xi)$ is uniquely determined by $w(p^{*}\xi)=p^{*}w(\xi)=w(L_{1})(\dots)w(L_{n})$.
\end{corollary}

\begin{remark}
This gives the following intuition: it explains why requiring
$w_{1}(\gamma_{1}^{1})=\widetilde{a}$ is enough for our axioms. Writing
$w(\gamma_{1}^{1})=1+\widetilde{a}$ where $\widetilde{a}\in H^{1}(\gamma_{1}^{1};\ZZ_{2})$ is
the generator, then the first thing we can do is pass to the universal
line bundle to prove $w(\gamma^{1})=1+a$. We know there exists a base
map $g_{1}\colon Y\to\RP^{\infty}$ such that if we pullback
$g_{1}^{*}\gamma^{1}\iso L_{1}$, then $g_{1}^{*}w(\gamma^{1})=w(L_{1})$.
We can do this for each $L_{1}$, \dots, $L_{n}$ to obtain $p$ uniquely.
This gives us a way to prove the uniqueness for Stiefel--Whitney classes.
\end{remark}

\begin{example}
Recall, we constructed a map
\begin{equation}
f\colon\RP^{\infty}\times\dots\times\RP^{\infty}\to G_{n}
\end{equation}
such that
\begin{enumerate}
\item the pullback on the cohomology ring $f^{*}\colon H^{*}(G_{n};\ZZ_{2})\to H^{*}((\RP^{\infty})^{n};\ZZ_{2})$
is injective, and
\item $f^{*}\gamma^{n}\iso\xi=\pi_{1}^{*}\gamma^{1}\oplus\dots\oplus\pi_{n}^{*}\gamma^{1}$
where $\pi_{j}\colon(\RP^{\infty})^{n}\to\RP^{\infty}$ is the
$j^{\text{th}}$ projection. [Here, confusingly, $\gamma^{n}$ plays the
  role of $\xi$ in the splitting principle, and more straightforwardly
  $\pi_{j}^{*}\gamma^{1}$ plays the role of $L_{j}$]
\end{enumerate}
\end{example}

\begin{definition}
Let $\xi$ be a (real) vector bundle $E\to B$.
The \define{Projectivization} of $\xi$ is the fiber bundle
$\RR\Projectivization{\xi}$ over $B$ where each fiber is just the
projectivization of $F_{b}$ in $\xi$.

Equivalently, $\RR\Projectivization{\xi}:=(E\setminus s_{0}(X))/\sim$
where $s_{0}$ is the zero-section, and $(x,v)\sim(x,\lambda v)$ for
$\lambda\in\RR\setminus\{0\}$. What is the fiber for such a thing?
Well, recall, $E$ has fiber $\RR^{n}$, then removing the zero section
gives $\RR^{n}\setminus\{0\}$, and then quotienting out by the
equivalence relation gives us $\RP^{n-1}$ as the fiber.

\emph{Note: this is no longer a vector bundle}. But it remains a fiber bundle.
\end{definition}



\begin{theorem}[Splitting principle --- reiterate]
Let $\xi\colon E\to X$ be an $\RR^{n}$-bundle, where $X$ is a
paracompact space. Then there exists a topological space $Y$ and a map
$p\colon Y\to X$ such that
\begin{enumerate}
\item $p^{*}\colon H^{*}(X;\ZZ_{2})\to H^{*}(Y;\ZZ_{2})$ is injective
\item the pullback bundle $p^{*}\xi$ splits
  $p^{*}E=L_{1}\oplus\dots\oplus L_{n}$ as the Whitney sum of line
  bundles $L_{1}$, \dots, $L_{n}$ over $Y$.
\end{enumerate}
\end{theorem}

\begin{goal}
Find a space $X'$ and a map $p\colon X'\to X$ such that
\begin{enumerate}
\item the induced morphism $p^{*}$ [on cohomology groups] is
  injective; and
\item $p^{*}E=p^{*}\xi=E'\oplus L$ where $L$ is a line bundle over
  $X'$ and $E'$ is a rank $(n-1)$-bundle over $X'$.
\end{enumerate}
\end{goal}

\begin{proof}[Proof (splitting principle)]
By induction on $n$.
Then since $n\in\NN$, we just iterate this procedure finitely many
times, and induction proves everything works fine.

We take $p\colon\RR\Projectivization{E}\to X$ which gives us
$p\colon X'\to X$ and satisfies $p^{*}\colon H^{n}(X;\ZZ_{2})\to H^{n}(X';\ZZ_{2})$
is injective (by Levy--Hirsch theorem). This achieves (1) of our goal.

Now we need to find the line bundle $L$ and the bundle $E'$. Then we
look at $p^{*}\xi$ which is a bundle over $X'$. A point in $X'$ is an
equivalence class $[(x,v)]$ and we want to find a subbundle of
$p^{*}\xi$ which is a line bundle.
We can look at $\Span(v)$ to get a line in $\RR^{n}$, then pulling it
back gives a line subbundle $L$ in $p^{*}\xi$. We use the fact that
$X$ is paracompact, so it has a Euclidean metric for $\xi$. Then we
can pullback this metric along $p^{*}$ to get a metric for
$p^{*}\xi$. Then we find the orthogonal complement $L^{\perp}$ such
that $L^{\perp}\oplus L\iso p^{*}\xi$.

Then we induct on $n$ to get the decomposition.
\end{proof}

\subsection{Existence of Stiefel--Whitney classes}

\begin{node}
We need two ingredients from algebraic topology:
\begin{enumerate}
\item Thom isomorphism
\item Steenrod squares
\end{enumerate}
\end{node}

\begin{notation}
Let $\xi\colon E\to B$ be an $\RR^{n}$-bundle. Let $E_{0}:=E\setminus s_{0}(B)$
be $E$ after removing its zero section. The fiber $F\iso\RR^{n}$,
$F_{0}\iso F\setminus\{0\}$. We have the inclusion map
\begin{equation}
i\colon(F,F_{0})\into(E,E_{0}),
\end{equation}
which induces the morphism on cohomology groups
\begin{equation}
i^{*}\colon H^{*}(E,E_{0};\ZZ_{2})\to H^{*}(F,F_{0};\ZZ_{2}),
\end{equation}
and we have the isomorphisms
\begin{equation}
H^{*}(F,F_{0};\ZZ_{2})\iso H^{*}(\RR^{n},\sphere{n};\ZZ_{2})\iso\widetilde{H}^{*}(\sphere{n};\ZZ_{2})
\end{equation}
where
\begin{equation}
\widetilde{H}^{k}(\sphere{n};\ZZ_{2})\iso\begin{cases}\ZZ_{2} &
\mbox{if }k=n\\
0 & \mbox{otherwise}
\end{cases}
\end{equation}
is the reduced cohomology of the sphere.
\end{notation}

\begin{theorem}[Thom]
\begin{enumerate}
\item There exists a unique $U\in H^{n}(E,E_{0};\ZZ_{2})$ such that
  $i^{*}U\in H^{n}(F,F_{0};\ZZ_{2})$ is the generator for
  $H^{n}(F,F_{0};\ZZ_{2})\iso\ZZ_{2}$ for every fiber $F$. We call $U$
  the \define{Thom Class}.
\item We have the isomorphism
  \begin{align}
     H^{i}(E;\ZZ_{2})&\to H^{i+n}(E,E_{0};\ZZ_{2})\\
 Y&\mapsto Y\smile U\nonumber
  \end{align}
  (obtained by mapping elements $Y$ to the cup product $Y\smile U$
  with $U$) which is called the \define{Thom Isomorphism}.
\end{enumerate}
\end{theorem}

\begin{remark}
We can write
\begin{align}
\varphi\colon & H^{i}(B;\ZZ_{2})\to H^{i+n}(E,E_{0};\ZZ_{2})\\
& \beta\mapsto p^{*}\beta\smile U\nonumber,
\end{align}
since $E$ deformation retracts to $B$ (since the fiber
$\RR^{n}\homotopic\point{*}$) via $p\colon E\to B$ we expect to be an isomorphism.
\end{remark}

\begin{node}[Foreshadowing]
How are Stiefel--Whitney classes defined using Thom isomorphism? We
have $w_{i}(\xi)\in H^{i}(B;\ZZ_{2})$. Define
$w_{i}(\xi):=\varphi^{-1}(\Sq_{i}(U))$ where $U\in H^{n}(E,E_{0};\ZZ_{2})$ is the 
Thom class, and
\begin{equation}
\Sq_{i}\colon H^{n}(E,E_{0};\ZZ_{2})\to H^{n+1}(E,E_{0};\ZZ_{2})
\end{equation}
is the Steenrod square, so
\begin{equation}
\varphi^{-1}\circ\Sq_{i}\colon H^{n}(E,E_{0};\ZZ_{2})\to H^{i}(B;\ZZ_{2}).
\end{equation}
Then we just need to prove this satisfies all the axioms.
\end{node}

\begin{node}[Proof sketch for Thom's isomorphism theorem]
Step 1: working in a local trivialization chart (or equivalently where
$E=B\times F$) prove everything works [requires K\"{u}nneth formula]. Step 2: consider the union of 2
local trivializing charts, and prove everything works [requires
  Mayer--Vietoris sequence]. Step 3: from
local to global (since $B$ is not assumed compact, nor even
paracompact, this requires a lot of work) [requires direct limit of cohomology].
\end{node}